\documentclass[a4paper,nobib,justified]{tufte-book}

\hypersetup{colorlinks}

%%
% If they're installed, use Bergamo and Chantilly from www.fontsite.com.
% They're clones of Bembo and Gill Sans, respectively.
%\IfFileExists{bergamo.sty}{\usepackage[osf]{bergamo}}{}% Bembo
%\IfFileExists{chantill.sty}{\usepackage{chantill}}{}% Gill Sans

\makeatletter

%%
% Just some sample text
%\usepackage{lipsum}

%%
% For nicely typeset tabular material
\usepackage{booktabs}
\usepackage{multirow}
\usepackage{adjustbox}
\usepackage{threeparttable}
\usepackage{array}

\newcolumntype{L}[1]{>{\noindent\RaggedRight\arraybackslash\hspace{0pt}}p{#1}}
\newcolumntype{R}[1]{>{\noindent\RaggedLeft\arraybackslash\hspace{0pt}}p{#1}}
\newcolumntype{C}[1]{>{\noindent\Centering\arraybackslash\hspace{0pt}}p{#1}}

%%
% For graphics / images
\usepackage{graphicx}
\setkeys{Gin}{width=\linewidth,totalheight=\textheight,keepaspectratio}
\graphicspath{{media/}}

% Add numbering to sections and show them in the TOC
\setcounter{secnumdepth}{3}
\setcounter{tocdepth}{3}

% Bibliography
\usepackage[
style=numeric,
citestyle=numeric-comp,
autocite=inline,
autopunct=true,
backend=biber,
maxbibnames=99,
maxcitenames=2,
mincitenames=1,
defernumbers=true,
urldate=iso,
]{biblatex}
\addbibresource{bibliography.bib}
\addbibresource{contributions.bib}
\DeclareBibliographyCategory{contributions}

%%
% Maths
\usepackage{amsmath}

% Multiple columns
\usepackage{multicol}

% Prints a trailing space in a smart way.
\usepackage{xspace}

% Boxes
\usepackage{tcolorbox}

% Inserts a blank page
\newcommand{\blankpage}{\newpage\hbox{}\thispagestyle{empty}\newpage}

% Referencing
\usepackage[unicode]{hyperref}
\usepackage[nameinlink]{cleveref}

\hypersetup{%
	bookmarksnumbered = true, % Show section numbering in the PDF table of contents. Allows easier browsing of the document
	colorlinks = true, % uncomment this line if you prefer colored hyperlinks (e.g., for onscreen viewing)
%	pdfborder = {0 0 0},
	bookmarksdepth = subsection,
%	citecolor = DarkGreen,
%	linkcolor = DarkBlue,
%	urlcolor = DarkGreen,
}

%\AtBeginDocument{%
%	\hypersetup{
%		pdfauthor={\plainauthor}
%		pdftitle={\plaintitle},
%	}%
%}

\usepackage{siunitx}
\usepackage[yyyymmdd]{datetime}
\usepackage{ifdraft}

% Draft packages
\usepackage{comment}

% Final packages
%\ifoptionfinal{
	%\usepackage[protrusion=true,expansion,babel=true]{microtype}
	\usepackage{microtype}
%}{}

% Modifications to the default tufte template to make it more applicable to a thesis
% Credit goes to Tiffany Tseng, lalider, see https://github.com/lalider/tufte-latex-thesis

\usepackage[parfill]{parskip}

% remove paragraph indentation

% Paragraph indentation and separation for normal text
\renewcommand{\@tufte@reset@par}{%
	\setlength{\RaggedRightParindent}{0.0pc}%
	\setlength{\JustifyingParindent}{0.0pc}%
	\setlength{\parindent}{0pc}%
	\setlength{\parskip}{\baselineskip}%
}
\@tufte@reset@par

% Paragraph indentation and separation for marginal text
\renewcommand{\@tufte@margin@par}{%
	\setlength{\RaggedRightParindent}{0.0pc}%
	\setlength{\JustifyingParindent}{0.0pc}%
	\setlength{\parindent}{0.0pc}%
	\setlength{\parskip}{10pt}%
}

\titleformat{\section}%
[hang]% shape
{\normalfont\Large}% format applied to label+text
{\thesection}% label
{1em}% horizontal separation between label and title body
{}% before the title body
[]% after the title body

%\titlespacing*{\section}{0pt}{3.5ex plus 1ex minus .2ex}{2.3ex plus .2ex}
%\titlespacing*{\subsection}{0pt}{3.25ex plus 1ex minus .2ex}{1.5ex plus.2ex}
\titlespacing*{\section}{0pt}{30pt}{20pt}
\titlespacing*{\subsection}{0pt}{20pt}{5pt}


% Disable adding empty pages without a purpose (not intended for printing in a book format...)
\def\cleardoublepage{
	\clearpage%
}

% Acronyms
\usepackage{acro}
\acsetup{
	use-id-as-short,
	make-links=false,
	case-sensitive=false,
	patch/floats=false,
	patch/caption=false,
	patch/tabularx=false
}

\NewAcroTemplate[list]{glossary}{%
	\begin{description}
		\acronymsmapF{%
			\item[\sffamily\textbf{\acrowrite{short}}] \acrowrite{long}%
			\acroifanyT {foreign,extra} {~(}%
			\acrowrite {foreign}%
			\acroifallT {foreign,extra} {,~}%
			\acrowrite {extra}%
			\acroifanyT {foreign,extra} {)}%
			\acropagefill%
			\acropages%
				{ \acrotranslate {page} \nobreakspace }%
				{ \acrotranslate {pages} \nobreakspace }%
		}%
		{ \item \AcroRerun{list} }%
		\end {description}
	}

% \acuse does not add page numbers to the acronym list. \acempty does.
\NewAcroTemplate{empty}{}
\NewAcroCommand\acempty{m}{\UseAcroTemplate{empty}{#1}}

% Various utilities
\usepackage{tabularx}
\usepackage[singlelinecheck=false,font={footnotesize}]{caption}
\usepackage{subcaption}
\usepackage{comment}
\usepackage{epigraph}


\usepackage{tikz}
\usepackage{pgfplots}
\pgfplotsset{compat = 1.18}
\usetikzlibrary{arrows.meta,positioning,automata,calc}
%\usepackage{physics}
\let\Re=\relax
\let\Im=\relax

% Colours
\definecolor{off}{HTML}{e53935}
\definecolor{on}{HTML}{43a047}
\usepackage{xcolor-solarized}
\usepackage{xcolor-material}
\usepackage{colortbl}

%\usepackage{todonotes}
\usepackage[inline]{enumitem}
\usepackage[absolute,overlay]{textpos}
\usepackage{pdflscape}
\usepackage{geometry}
\usepackage{float}

% Code
\usepackage{minted}
\usepackage{xpatch,letltxmacro}
\LetLtxMacro{\cminted}{\minted}
\let\endcminted\endminted
\xpretocmd{\cminted}{\RecustomVerbatimEnvironment{Verbatim}{BVerbatim}{}}{}{}
%\definecolor{mintedbg}{rgb}{0.95,0.95,0.93}
\setminted{
	% bgcolor=mintedbg,
	tabsize=2
}
\setminted[text]{baselinestretch=0.8}

% Git magic
\write18{git log --pretty='\@percentchar Creset\@percentchar s' --no-merges -1 > \jobname.git1.tmp}
\write18{git rev-parse --short HEAD > \jobname.git2.tmp}
\def\gitcommit{\input{\jobname.git2.tmp}\unskip}
\def\gitcommitmessage{\input{\jobname.git1.tmp}\unskip}

% Smart citations

\xdef\@NextCitationOffsetTemp{0pt}%
\NewDocumentCommand{\NextCitationOffset}{m}{%
	\xdef\@NextCitationOffsetTemp{#1}%
}%

\newbibmacro*{cite:title}{%
  \printtext[bibhyperref]{%
    \printfield[citetitle]{labeltitle}}%
  }

\newbibmacro*{cite:shorthand}{%
	\printtext[bibhyperref]{\printfield{shorthand}}}

\newbibmacro*{cite:full}{%
	\iffieldundef{shorthand}
	{%
		\printnames{labelname}%
		\iffieldundef{year}{}{\setunit*{\addspace}\printtext[parens]{\printfield{year}}}%
		\setunit*{\printdelim{nametitledelim}}%
		\printtext[bibhyperref]{%
    	\printfield[citetitle]{labeltitle}}%
	}%
	{\usebibmacro{cite:shorthand}}}

\newbibmacro*{cite:marginnote}{%
	\printtext{%
		\marginnote[\@NextCitationOffsetTemp]{%
			\printtext{[\printfield{labelprefix}\printfield{labelnumber}]}%
			\setunit*{\addspace}%
			\usebibmacro{cite:full}%
			\finentry%
		}%
	}%
	\xdef\@NextCitationOffsetTemp{0pt}%
}


\DeclareCiteCommand{\dualcite}
  {\usebibmacro{cite:init}%
   \usebibmacro{prenote}}%
  {\usebibmacro{citeindex}%
   \usebibmacro{cite:marginnote}%
   \usebibmacro{cite:comp}%
   }%
  {}%
  {\usebibmacro{cite:dump}%
   \usebibmacro{postnote}%
   }

\DeclareMultiCiteCommand{\dualcites}[\mkbibbrackets]{\dualcite}{\multicitedelim}

\def\tuftecite{
	\iffootnote{footnote}{text}
}

% Patching commands that enter information in margins, so that we can detect which
% citation should be used
\newtoggle{inmargin}
\pretocmd{\marginpar}{\toggletrue{inmargin}}{}{}
\apptocmd{\@xympar}{\togglefalse{inmargin}}{}{}
\pretocmd{\@tufte@float}{\toggletrue{inmargin}}{}{}
\pretocmd{\end@tufte@float}{\togglefalse{inmargin}}{}{}
\pretocmd{\@tufte@margin@float}{\toggletrue{inmargin}}{}{}
\pretocmd{\end@tufte@margin@float}{\togglefalse{inmargin}}{}{}

\def\autocite{%
	\iftoggle{inmargin}\parencite\dualcites%
}

% Tufte disables subparagraphs. This brings them back.
\renewcommand\subparagraph{\@startsection{subparagraph}{5}{\parindent}%
	{3.25ex \@plus1ex \@minus .2ex}%
	{-1em}%
	{\hspace{1em}\normalfont\normalsize\itshape}}

% Definitions
\usepackage{keytheorems}

\newsavebox{\definition@caption@box}
\newsavebox{\definition@content@box}
\newsavebox{\definition@margin@box}

\newkeytheoremstyle{margincaptionstyle}{
	headpunct={\hspace{0pt}},
	postheadspace={0pt},
	preheadhook = {%
		\def\deferred@thm@head#1{%
			\everypar\dth@everypar%
			\sbox\@labels{}%
			\global\sbox{\definition@caption@box}{%
				\begin{minipage}[t]{\marginparwidth}%
					\vspace{\baselineskip}%
					\normalfont #1%
				\end{minipage}%
			}%
		}%
		\begin{lrbox}{\definition@content@box}%
		\begin{minipage}[t]{\textwidth}%
			\@tufte@reset@par%
	},
	postfoothook = {%
		\vspace{-\baselineskip}%
		\end{minipage}%
		\end{lrbox}%
		%
		\begin{fullwidth}%
			\@tufte@float@textwidth[\@tufte@caption@vertical@offset]{\definition@content@box}{\definition@caption@box}%
		\end{fullwidth}%
	}
}

\newkeytheoremstyle{marginnotestyle}{
	headpunct = {:},
	preheadhook = {%
		\begin{lrbox}{\definition@margin@box}%
		\begin{minipage}[t]{\marginparwidth}%
			\@tufte@marginnote@font\@tufte@marginnote@justification\@tufte@margin@par\vspace*{-1\baselineskip}%
	},
	postfoothook = {%
		\end{minipage}%
		\end{lrbox}%
		\marginnote{\usebox\definition@margin@box}%
	}
}

\newkeytheorem{definition}[
	name=Definition,
	parent=chapter,
	style=margincaptionstyle
]
\newkeytheorem{margindefinition}[
	name=Definition,
	sibling=definition,
	style=marginnotestyle,
	refname=Definition
]

% \acuse command that also adds the current page
%\NewAcroTemplate{fdirthesisempty}{}
%\NewAcroCommand\acusepage{m}{}
\let\acusepage\acuse

% Acronym on margin command
\usepackage{suffix}
\NewAcroTemplate{margin-format}{%
  \emph{\acrowrite{short}:} \acrowrite{long}%
  \acroifT{extra}{, \acrowrite{extra}}%
}
\NewAcroCommand\marginacrotext{m}{\UseAcroTemplate{margin-format}{#1}}
\NewDocumentCommand{\mac}{ O{0pt} m }{%
  \marginnote[#1]{\vspace{1ex}\marginacrotext*{#2}\vspace{1ex}}%
}
\WithSuffix\NewDocumentCommand\acs+{ O{0pt} m }{%
  \acs{#2}\mac[#1]{#2}%
}

% Let us have more floats
\extrafloats{100}

\definecolor{draft}{HTML}{bf616a}

\NewDocumentCommand\draft{m}{%
	\textcolor{draft!65!black}{#1}%
}

\NewDocumentCommand\todo{m}{%
	\textcolor{draft}{#1}%
}

\makeatletter

%\newcommand\draft{\bgroup\color{draft!50!white}\markoverwith
%	{\textcolor{black}{\rule[-0.5ex]{2pt}{0.4pt}}}\ULon}

\newcommand\footurl@[1]{\footnote{\url@{#1}}}
\DeclareRobustCommand{\footurl}{\hyper@normalise\footurl@}

\newcommand\foothref@[2]{\href@{#1}{#2}\hyper@punct\footnote{\url@{#1}}}
\def\Hy@foothref#{%
	\hyper@normalise\foothref@
}
\DeclareRobustCommand*{\foothref}[1][]{%
	\begingroup%
	\global\def\hyper@punct{#1}%
	\@ifnextchar\bgroup\Hy@foothref{\hyper@normalise\foothref@}%
}
\makeatother

\newcommand\numberthis{\addtocounter{equation}{1}\tag{\theequation}}

% Transitions
\AtBeginDocument{
\colorlet{invalid}{MaterialAmber900}
\colorlet{unexpected}{MaterialRed900}
}

\def\ar{ \( \rightarrow \) }

\makeatletter

\makeatother

\makeatother

%%
% Book metadata
\title{Systems Engineering for Sub-CubeSat Spacecraft\thanks{AcubeSAT}}
\author[Konstantinos Kanavouras]{Konstantinos Kanavouras}
\publisher{University of Luxembourg}

\acsetup{foreign/display=false, list/foreign/display=false}

\def\g{\ignorespaces}
\def\e{\ignorespaces}
\def\acusepage#1{}

\usepackage{microtype}

\hypersetup{
	pdftitle={Systems Engineering for Sub-CubeSat Spacecraft},
	pdfsubject={PhD Thesis},
	pdfauthor={Konstantinos Kanavouras},
	addtopdfcreator={tufte-book class}
}

\begin{document}
	
\addcontentsline{toc}{chapter}{Preamble}

% Front matter
%\frontmatter

% r.1 blank page
%\blankpage

% r.3 full title page
\makeatletter
\renewcommand{\maketitle}{%
	\newpage
	\global\@topnum\z@% prevent floats from being placed at the top of the page
	\begingroup
	\setlength{\parindent}{0pt}%
	\setlength{\parskip}{4pt}%
	\let\@@title\@empty
	\let\@@author\@empty
	\let\@@date\@empty
	\thispagestyle{empty}
	\begin{fullwidth}
		\vfill
		\begin{center}
			University of Luxembourg Logo\par
			\vspace{1cm}
			\LARGE\textsc{Doctoral Thesis}\par
			\vspace{6ex}
			\hrule
			\vspace{4ex}
			\Huge\textbf{Systems Engineering}\\[1ex]
			\LARGE\textbf{for}\\
			\Huge\textbf{Sub-CubeSat Spacecraft}\par
			\vspace{2.7ex}
			\hrule

			\vspace{6ex}

			\Large
			\begin{tabular}{ll}
				\emph{Author:} & \href{https://github.com/kongr45gpen}{Konstantinos \textsc{Kanavouras} \normalsize (\fontfamily{pplx}\selectfont 8824)}
				\\[1.5ex]
				\emph{Supervisor:} & \href{http://ee.auth.gr/en/school/faculty-staff/electronics-computers-department/hatzopoulos-alkiviadis/}{Prof. Andreas M.~\textsc{Hein}}
			\end{tabular}

			\vspace{6ex}

			\large \textit{A thesis submitted in fulfillment of the requirements\\ for something}\\[0.3cm] % University requirement text
			\textit{in the}\\[0.4cm]
			\href{https://www.eng.auth.gr/en/home.html}{Faculty of Engineering}
			\\
			\href{https://ee.auth.gr/en/}{Electrical \& Computer Engineering Department}
			\\[1cm] % Research group name and department name

			\vspace{6ex}

			{\large June 28, 2021}\\[4cm] % Date

		\end{center}
		\vfill
	\end{fullwidth}
	\endgroup
	\thispagestyle{plain}% suppress the running head
	\tuftebreak% add some space before the text begins
	\@afterindentfalse\@afterheading% suppress indentation of the next paragraph
}
\makeatother
\maketitle

% v.4 copyright page
\newpage
\phantomsection
\addcontentsline{toc}{section}{Copyrights}
\begin{fullwidth}
~\vfill
\thispagestyle{empty}
\setlength{\parindent}{0pt}
\setlength{\parskip}{\baselineskip}
Copyright \copyright\ \the\year\ \thanklessauthor

%\par\smallcaps{\href{https://github.com/kongr45gpen/phd-thesis}{Https://Github.com/Kongr45gpen/PhD-Thesis}}

\par Licensed under the Creative Commons Attribution 4.0 International (CC BY 4.0) License (the ``License''); you may not
use this file except in compliance with the License. You may obtain a copy
of the License at \url{https://creativecommons.org/licenses/by/4.0/legalcode}, or a human-readable summary at \url{https://creativecommons.org/licenses/by/4.0/}.
Unless
required by applicable law or agreed to in writing, material distributed
under the License is distributed on an \smallcaps{``AS IS'' BASIS, WITHOUT
WARRANTIES OR CONDITIONS OF ANY KIND}, either express or implied. See the
License for the specific language governing permissions and limitations
under the License.\index{license}

\par This thesis was produced using \LaTeX, the \href{https://ctan.org/pkg/tufte-latex?lang=en}{\texttt{tufte-latex}} template, and the modifications from \linebreak[4] \href{https://github.com/lalider/tufte-latex-thesis}{\texttt{tufte-latex-thesis}}.

\par\textit{Rendered on \today{} (\texttt{\gitcommit}: \gitcommitmessage)}
\end{fullwidth}

% r.5 contents
\tableofcontents
\addcontentsline{toc}{section}{Contents}

\begin{fullwidth}
\listoffigures
\addcontentsline{toc}{section}{List of Figures}

\listoftables
\addcontentsline{toc}{section}{List of Tables}

%\chapter*{List of Acronyms}
%\acuseall%
\bgroup
\setlength\parskip{1ex}
\printacronyms[pages={display=all,seq/use=false},]
\addcontentsline{toc}{section}{Acronyms}
\egroup

\end{fullwidth}

% r.9 introduction
\cleardoublepage

\chapter*{Abstract}

\justify
Space is not a welcoming environment.

%%
% Start the main matter (normal chapters)
\mainmatter


\chapter{Introduction}

\section{Thesis research area}


\begin{marginfigure}
	\centering
	(Imagine a chart here)
	\caption[Types of satellites in operation in 2018]{Types of satellites in operation in 2018 \autocite{ECSS-E-ST-10-02C}}
	\label{fig:satellite_types}
\end{marginfigure}

This thesis involves satellite design and operations, parts of the Aerospace Engineering disciplines. The first artificial satellite \e{``Sputnik 1''} was launched in 1957, and by 2021 about 8900 satellites have been launched, of which about 4000 are currently operational and in orbit \autocite{ECSS-E-ST-10-02C}.

Artificial satellites can perform a number of different functions, which typically include telecommunications, earth observation, navigation (\acs{GNSS}), scientific research, technology demonstration, and other purposes, excluding manned missions, missions beyond Earth orbit, and space stations. The data generated using space technology have long been an integral part of modern civilisation, having influenced transportation, weather forecasting, and technology in general.

\begin{figure}
	\centering
    (satellites here)
	\caption[Active satellites orbiting the earth in 2021]{Active satellites orbiting the earth in 2021. \acl{LEO} satellites can be seen close to the planet, and Geostationary Orbit satellites can be seen at a fixed high altitude, while Medium Earth Orbit are between the two previous categories.}
    \label{fig:all_satellites}
    \setfloatalignment{t}
\end{figure}
			
After launch, satellites typically remain in space until the completion of their mission after some years, at which point they can either return to Earth and burn up in the atmosphere, or remain in space as ``space junk''.

\appendix

\begin{fullwidth}
\printbibliography[heading=bibnumbered]
\end{fullwidth}


\printindex



\end{document}

